\documentclass[a4paper,12pt]{article}
\usepackage[margin=1in]{geometry}
\usepackage{enumitem}
\usepackage{parskip}

\begin{document}

% --- Header ---
\begin{center}
    \textbf{\Large RV College of Engineering} \\
    Bangalore \\
    \textbf{IP Cell}
\end{center}

\vspace{0.5cm}

\section*{Invention Disclosure Guidelines}

Following guidelines are provided in order to channelize your ideas, project work, concepts, invention and bring out the creative work done.

\begin{itemize}[itemsep=0.5em]
    \item You need document your invention in such a way that as to be very clear to be one who is not familiar with the work.
    \item Describe, what makes your invention different from what has done before.
    \item In addition to describing all the modules, functions, parts, describe how they work together.
    \item Why did you do things the way you did them, and not some other way?
    \item How else could you have accomplished the same end?
    \item Do not limit yourself to exactly the idea can be built, or prototype or proof-of-concept you have with you.
    \item Open for your imagination to run - how else this invention work?
    \item What are the possible problems? Under what circumstances might your invention not work?
    \item While preparing invention disclosure form, prepare drawings, figures, flow charts so that it will be easy to understand.
    \item Disclose your inventions as early as possible to IP Cell in order to protect then if patentable.
    \item Do not disclose your inventions to anyone in any form until it is protected in the form of IP.
    \item Do not publish the invention in any form like presentations in seminar, conferences, journals, newspaper, magazines, customer meetings, TV shows, any media and so on.
    \item Inventions that are generated as part of company sponsored projects or DST funded projects are to be informed clearly and this can be filed as joint IP if they agree.
    \item Report inventions that are disruptive, have commercial value and useful to mankind.
    \item Your invention will qualify for patenting only if it meets patentability conditions – novelty, inventive step and have industrial application.
\end{itemize}

\end{document}